\section{Conclusion}
\label{sec:conclusion}

This paper presented a systematic analysis of noise robustness in encoder-only transformer architectures, identifying critical vulnerability transitions at layers 3 and 8 that correspond to boundaries between linguistic processing phases. Through evaluation of 300,000 perturbed samples across five models, we demonstrated that these transitions represent universal computational properties, with 61.1\% correlation in vulnerability patterns across architectures.

Our key findings include: (1) RoBERTa's superior robustness (98.8\% average) stems from training choices that align with natural phase boundaries, particularly dynamic masking and larger batch sizes; (2) Character-level perturbations show 85\% recovery through semantic layers while syntactic disruptions cause 78\% degradation, reflecting the brittleness of hierarchical syntactic processing; (3) Strategic layer dropout at identified transitions achieves theoretical 3.1× speedup while maintaining 95\% performance, though actual runtime validation remains necessary.

The theoretical analysis suggests these transitions emerge from information-theoretic optimization during training, where models naturally develop distinct phases for surface, syntactic, and semantic processing. This aligns with linguistic theory while providing a computational framework for understanding transformer vulnerability.

\textbf{Limitations:} Our analysis is restricted to encoder architectures and English text. Decoder models (GPT, LLaMA) may exhibit different patterns aligned with generation rather than understanding. The claimed speedups are theoretical calculations requiring empirical validation in production environments.

\textbf{Practical Implications:} For deployment in noise-critical applications, we recommend RoBERTa-based architectures, implementation of quality-aware routing for adaptive processing, and targeted denoising at identified vulnerable layers. These strategies can reduce computational costs while maintaining robustness.

\textbf{Future Directions:} Important research areas include: (1) Systematic analysis of decoder architectures to identify generation-specific vulnerabilities; (2) Multilingual studies to determine universality of transitions; (3) Development of phase-aware architectures that explicitly model transition boundaries; (4) Runtime validation of theoretical efficiency gains in production systems.

The identification of universal phase transitions advances our understanding of transformer architectures beyond black-box models toward interpretable systems with predictable vulnerability patterns. As transformer models become increasingly critical in real-world applications, this knowledge enables development of more robust and efficient NLP systems that can reliably handle the noisy, imperfect data encountered in practice.